% GENERATED FILE. Edit canonical lessons, then run npm run generate:books.
% canonical-source-hash: fnv1a64:cb8a594e1def2dcd
% canonical-lessons: MW-R26-shopping-counter, MW-R26-transport-counter, MW-R26-food-counter, MW-R26-letters-close, MW-C26-counter-exchange, MW-R26-script-close

\chapter{One Exchange, Three Counters}
\label{ch:mw-counter-exchange}
\hlchaptermodality{31}

\begin{chapteropening}
\textbf{By the end of this chapter:} \emph{I can run one five-turn counter exchange over the shopping, transport, and food words this book has already taught, changing nothing but the noun.}

Each of the three noun sets is run through the exchange in turn, the five least-drilled signs are retrieved alone, and the whole five-turn exchange then passes separately in listening, speaking, reading, and writing.
\end{chapteropening}

\section[Practice]{The shopping map, inside the exchange}

\label{lesson:MW-R26-shopping-counter}



\begin{quote}
Recall all three shopping payoffs in the order they were earned: three words, five words, seven words.
\end{quote}

\subsection*{Your turn}
Take the seven shopping words one at a time. For each one, say its meaning, then put it into the request from Chapter 26 and the price question from Chapter 27. Two of the seven do not fit a request slot at all: name them and say why.

Then write all seven from dictation, and write the two lines you built with the word that fits them best.

\begin{tcolorbox}[breakable,colback=teal!4,colframe=teal!35!black,title={Before you move on}]
Source: \href{https://www.marwaripathshala.com/marwari-lesson-9-english}{Marwari Pathshala, Lesson 9}.
\end{tcolorbox}
\section[Practice]{The transport map, inside the exchange}

\label{lesson:MW-R26-transport-counter}



\begin{quote}
Recall both transport payoffs, say lentils, then write the total question.
\end{quote}

\subsection*{Your turn}
Take the five transport words one at a time. Say each meaning, then price it with the question from Chapter 27 and answer with the counter-offer from Chapter 28. A fare is bargained exactly like a price, which is why no new line is needed here.

Then write all five from dictation, and write the price question with the word you would actually use standing at a roadside.

\begin{tcolorbox}[breakable,colback=teal!4,colframe=teal!35!black,title={Before you move on}]
Sources: \href{https://www.marwaripathshala.com/marwari-lesson-10-english}{Marwari Pathshala, Lesson 10}; \href{https://www.marwaripathshala.com/marwari-lesson-9-english}{Marwari Pathshala, Lesson 9}.
\end{tcolorbox}
\section[Practice]{The food map, inside the ordering line}

\label{lesson:MW-R26-food-counter}



\begin{quote}
Recall all three food payoffs in the order they were earned, then write \textbf{\mw{ट}} and \textbf{\mw{य}}.
\end{quote}

\subsection*{Your turn}
Take the seven food words one at a time. Say each meaning, then order it with the line from Chapter 30 and price it with the question from Chapter 27. Only the first word of the order changes; the asking word never does.

Then write all seven from dictation, and write three complete orders: the one you would ask for first, the one you would ask for last, and the drink.

\begin{tcolorbox}[breakable,colback=teal!4,colframe=teal!35!black,title={Before you move on}]
Source: \href{https://www.marwaripathshala.com/marwari-lesson-6-english}{Marwari Pathshala, Lesson 6}.
\end{tcolorbox}
\section[Practice]{Five signs, on their own}

\label{lesson:MW-R26-letters-close}



\begin{quote}
Recall the six-word weather payoff, then write weather, heat, and cold.
\end{quote}

\subsection*{Your turn}
Some signs in this track have been read a hundred times inside words and asked for on their own only once. This lesson asks for them alone, which is the only way to know whether they were learned.

Write each from a spoken cue, with nothing else on the line: \textbf{\mw{ख}}, \textbf{\mw{ट}}, \textbf{\mw{श}}, \textbf{\mw{ड़}}, \textbf{\mw{ौ}}, \textbf{\mw{ग}}, \textbf{\mw{ड}}, \textbf{\mw{ु}}, \textbf{\mw{ई}}. Then write one word that carries each: \textbf{\mw{आखरी}}, \textbf{\mw{रोटी}}, \textbf{\mw{रिक्शा}}, \textbf{\mw{थोड़ु}}, \textbf{\mw{मौसम}}, \textbf{\mw{गर्मी}}, \textbf{\mw{ठंडी}}, \textbf{\mw{दुकान}}, \textbf{\mw{घी}}. Repair only the missed sign and rewrite its word once.

\begin{tcolorbox}[breakable,colback=teal!4,colframe=teal!35!black,title={Before you move on}]
Source: \href{https://www.unicode.org/charts/PDF/U0900.pdf}{Unicode Devanagari chart}.
\end{tcolorbox}
\section[Practice]{Payoff — one exchange, three counters}

\label{lesson:MW-C26-counter-exchange}



\begin{quote}
Recall the five payoffs of this slice in order, then write clarified butter.
\end{quote}

\subsection*{Your turn}
Five turns, and the only thing that changes between counters is the noun.

\begin{enumerate}
\raggedright
  \item Hear the five turns shuffled and put them in order, twice: once with a shopping noun, once with a transport noun.
  \item Produce the whole exchange from meaning cues at a food stall, using the ordering line in place of the first turn.
  \item Read all five printed turns and say who speaks each.
  \item Write the whole exchange from dictation, without a model.
\end{enumerate}

Pass each skill separately. This is one function taught once and reused over three noun sets, not three exchanges learned three times. It still stops where numbers begin: no amount a seller names can yet be understood, and refusing the final price and leaving is untaught.

\begin{tcolorbox}[breakable,colback=teal!4,colframe=teal!35!black,title={Before you move on}]
Sources: \href{https://www.marwaripathshala.com/marwari-lesson-9-english}{Marwari Pathshala, Lesson 9}; \href{https://www.marwaripathshala.com/marwari-lesson-6-english}{Marwari Pathshala, Lesson 6}; \href{https://www.marwaripathshala.com/marwari-lesson-10-english}{Marwari Pathshala, Lesson 10}.
\end{tcolorbox}
\section[Practice]{Close the writing loop for \mw{ख}}

\label{lesson:MW-R26-script-close}



\begin{quote}
Recall the five-turn exchange payoff, then write vegetables and clarified butter.
\end{quote}

\subsection*{Your turn}
Write \textbf{\mw{ख}} and \textbf{\mw{क}} from breath cues alone, with nothing else on the line. Then write four words from dictation and mark which of the two each one uses: \textbf{\mw{दिखावो}}, \textbf{\mw{आखरी}}, \textbf{\mw{करो}}, \textbf{\mw{कितणो}}. Two take the breath and two do not.

Finish by writing every word this slice added, from meaning cues rather than sounds: \textbf{\mw{ये}}, \textbf{\mw{लो}}, \textbf{\mw{दो}}, \textbf{\mw{लावो}}, \textbf{\mw{घणो}}, \textbf{\mw{थोड़ु}}, and the five turns of the exchange. Repair only what was missed and rewrite that alone.

\begin{tcolorbox}[breakable,colback=teal!4,colframe=teal!35!black,title={Before you move on}]
Sources: \href{https://www.marwaripathshala.com/marwari-lesson-9-english}{Marwari Pathshala, Lesson 9}; \href{https://www.unicode.org/charts/PDF/U0900.pdf}{Unicode Devanagari chart}.
\end{tcolorbox}
